\chapter{多项式逼近}
\label{chap:polynomial-approximation-solutions}

\begin{solution}{bernstein-moment-identities}
由 $\sum_k\binom nkx^k(1-x)^{n-k}=1$ 得第一式。恒等式
$k\binom nk=n\binom{n-1}{k-1}$ 给出
\[
 \sum_k\frac kn\binom nkx^k(1-x)^{n-k}
 =x\sum_{j=0}^{n-1}\binom{n-1}jx^j(1-x)^{n-1-j}=x.
\]
当 $n\ge2$ 时，对 $2\le k\le n$ 使用
$k(k-1)\binom nk=n(n-1)\binom{n-2}{k-2}$，求和得到
\[
 \sum_{k=2}^n k(k-1)\binom nkx^k(1-x)^{n-k}=n(n-1)x^2.
\]
再用 $k^2=k(k-1)+k$ 得
$B_n(t^2)(x)=((n-1)/n)x^2+x/n$。$n=1$ 时直接算得 $B_1(t^2)(x)=x$，
仍满足该式。因此对所有 $n\ge1$，
$B_n((t-x)^2)=B_nt^2-2xB_nt+x^2B_n1=x(1-x)/n$。
\end{solution}

\begin{solution}{bernstein-positive-contraction}
$B_n$ 对采样值线性，故是线性算子。若 $f\ge0$，所有系数和基函数非负，
所以 $B_nf\ge0$。权重之和为 $1$，故 $B_n1=1$，并且
\[
 |B_nf(x)|\le B_n|f|(x)\le\norm f_\infty B_n1(x)=\norm f_\infty.
\]
取上确界得到所需范数估计；常函数表明算子范数等于 $1$，并非严格压缩。
\end{solution}

\begin{solution}{bernstein-modulus-estimate}
把 $f(k/n)-f(x)$ 乘以非负权重并求和，先用
\eqref{eq:modulus-telescoping} 得
\[
 |B_nf(x)-f(x)|\le\omega_f(\delta)
 \left(1+\delta^{-1}\sum_k|k/n-x|p_{n,k}(x)\right).
\]
加权 Cauchy--Schwarz 不等式和权重和为一给出
\[
 \sum_k|k/n-x|p_{n,k}
 \le\left(\sum_k(k/n-x)^2p_{n,k}\right)^{1/2}
 =\sqrt{x(1-x)/n}\le1/(2\sqrt n).
\]
取 $x$ 的上确界即得定理估计；令 $\delta=n^{-1/4}$，两个因子分别满足
$\omega_f(\delta)\to0$ 和括号有界。
\end{solution}

\begin{solution}{bernstein-lipschitz-rate}
由 Lipschitz 条件，
\[
 |B_nf(x)-f(x)|\le L\sum_k|k/n-x|p_{n,k}(x)
 \le L\sqrt{x(1-x)/n}\le\frac{L}{2\sqrt n}.
\]
估计与 $x$ 无关，故可取上确界。
\end{solution}

\begin{solution}{bernstein-linear-quadratic}
矩恒等式给出
\[
 B_n1=1,\qquad B_nt=t,
 \qquad B_nt^2=t^2+\frac{t(1-t)}n.
\]
因此所有仿射函数均被保持。固定 $n\ge1$。若 $p$ 的次数 $d>n$，
则 $B_np$ 次数至多为 $n$，不能等于 $p$。
以下排除 $2\le d\le n$。记下降阶乘
\[
 (k)_r=k(k-1)\cdots(k-r+1),\qquad (k)_0=1.
\]
多项式 $(k)_r$ 的次数为 $r$、首项系数为 $1$，故逐次消去最高次项可写成
$k^d=(k)_d+\sum_{r=0}^{d-1}c_r(k)_r$。对 $0\le r\le n$，组合恒等式给出
\[
 \sum_{k=0}^n(k)_r\binom nkx^k(1-x)^{n-k}=(n)_r x^r.
\]
这里 $k<r$ 的下降阶乘为零；$k\ge r$ 时把
$(k)_r\binom nk=(n)_r\binom{n-r}{k-r}$ 代入，再用二项式定理。
因此 $B_n(t^d)$ 的 $x^d$ 系数为
\[
 \frac{(n)_d}{n^d}=\prod_{j=0}^{d-1}\left(1-\frac jn\right)<1
 \qquad(2\le d\le n),
\]
而较低次单项式经 $B_n$ 后次数不会升高。若 $p$ 的首项系数为 $a_d\ne0$，
则 $B_np$ 的首项系数为 $a_d(n)_d/n^d\ne a_d$，所以仍不能等于 $p$。
结论是：对每个固定 $n\ge1$，所有不动多项式恰为 $p(t)=a+bt$，包括零多项式。
\end{solution}

\begin{solution}{bernstein-monotonicity}
利用组合恒等式整理导数可得
\[
 (B_nf)'(x)=n\sum_{k=0}^{n-1}
 \left[f\left(\frac{k+1}{n}\right)-f\left(\frac{k}{n}\right)\right]
 \binom{n-1}k x^k(1-x)^{n-1-k}.
\]
若 $f$ 单调递增，方括号非负，基函数也非负，故导数非负；递减情形同理。
\end{solution}

\begin{solution}{fixed-degree-polynomials-closed}
若 $a=b$，所有限制函数均为单点上的常函数，所得空间等于一维空间
$C(\{a\})$，因而闭。若 $a<b$，$1,x,\ldots,x^N$ 在区间上线性无关：
非零多项式不可能在整个区间上为零。因此限制空间维数为 $N+1$。
前文已经证明有限维赋范子空间
中的范数与坐标欧氏范数等价，因而它在任意赋范空间中完备。完备子空间必闭：
若其序列在母空间中收敛，则是 Cauchy 列，极限由子空间完备性仍落在其中。
故该多项式空间闭。
\end{solution}

\begin{solution}{dense-subspace-propagates-linear-identity}
设连续线性泛函 $L$ 在稠密集 $D$ 上为零。任给 $x$，取 $x_n\in D$ 使
$x_n\to x$。连续性给出 $L(x)=\lim_nL(x_n)=0$。所以 $L$ 在全空间为零。
此处实际只用了 $L$ 的连续性；线性结构不是这一步极限论证所必需的。
\end{solution}

\begin{solution}{absolute-value-polynomial-approximation}
若 $M=0$，可直接取 $p_n=0$。以下设 $M>0$。
若 $q_n\to|t|$ 于 $[-M,M]$，令 $p_n(t)=q_n(t)-q_n(0)$，则
$p_n(0)=0$ 且
$\norm{p_n-|\cdot|}_\infty\le
\norm{q_n-|\cdot|}_\infty+|q_n(0)|\to0$。若
$\norm f_\infty\le M$，便有
$\norm{p_n\circ f-|f|}_\infty$ 不超过同一个标量区间上的逼近误差。由于
$p_n$ 无常数项时仍由 $f$ 的有限次乘法和线性组合构成，它属于相应代数闭包。
\end{solution}

\begin{solution}{two-point-interpolation-from-separation}
由分离点性取 $h\in A$ 使 $h(x)\ne h(y)$。令
\[
 g= a+\frac{b-a}{h(y)-h(x)}(h-h(x)).
\]
含常数和代数的线性封闭性保证 $g\in A$，直接代入得 $g(x)=a$、$g(y)=b$。
\end{solution}

\begin{solution}{stone-weierstrass-finite-cover-proof}
空集情形平凡，以下设 $K\ne\varnothing$。给定 $f\in C(K;\R)$ 和
$\varepsilon>0$，由两点插值为每对 $x,y$ 取 $g_{x,y}\in A$，使其在
$x,y$ 处与 $f$ 相等；$x=y$ 时取常函数。固定 $x$，开集
\[
 V_y=\{z\in K:g_{x,y}(z)>f(z)-\varepsilon/2\}
\]
包含 $y$，因而覆盖 $K$。选有限子覆盖 $V_{y_1},\ldots,V_{y_r}$，令
$g_x=\max_{1\le j\le r}g_{x,y_j}$。格闭包引理保证 $g_x\in\overline A$；
每一点至少落在一个 $V_{y_j}$ 中，故 $g_x>f-\varepsilon/2$ 处处成立。
又因每个被取最大值的函数在 $x$ 处都等于 $f(x)$，有 $g_x(x)=f(x)$。

于是 $U_x=\{z:g_x(z)<f(z)+\varepsilon/2\}$ 是 $x$ 的开邻域。
再取有限子覆盖 $U_{x_1},\ldots,U_{x_s}$，令
$g=\min_{1\le i\le s}g_{x_i}\in\overline A$。
每个 $g_{x_i}$ 都有上述全局下界；在任一点，又至少有一个 $g_{x_i}$
满足相应上界。因此
\[
 f-\varepsilon/2<g<f+\varepsilon/2\quad\text{于 }K,
 \qquad \norm{g-f}_\infty\le\varepsilon/2.
\]
最后由 $g\in\overline A$ 取 $h\in A$ 使
$\norm{h-g}_\infty<\varepsilon/2$，得到 $\norm{h-f}_\infty<\varepsilon$。
第一次取最大值保留各处的下界，第二次取最小值补上各处的上界；有限性使
两次运算均留在代数闭包中。
\end{solution}

\begin{solution}{multivariate-polynomial-density}
若 $K$ 为空，结论平凡。否则多元多项式的限制组成 $C(K;\R)$ 的子代数并含常数。若 $x\ne y$，至少有
一个坐标 $j$ 满足 $x_j\ne y_j$，坐标函数 $p(z)=z_j$ 便分离这两点。实
Stone--Weierstrass 定理于是给出该代数在 $C(K;\R)$ 中一致稠密。
\end{solution}

\begin{solution}{trigonometric-polynomial-density}
把周期函数看成单位圆周 $\mathbb T$ 上的连续函数。有限和
$\sum_{|k|\le N}c_kz^k$ 构成复子代数，含常数；坐标函数 $z$ 分离圆周点；
共轭把 $z^k$ 变成 $z^{-k}$，仍在代数中。复 Stone--Weierstrass 定理给出
一致稠密性。写 $z=e^{it}$ 即得到通常的三角多项式。
\end{solution}

\begin{solution}{even-polynomial-density}
偶多项式的一致极限仍是偶函数，所以闭包包含于
\[
 E=\{f\in C([-1,1]):f(x)=f(-x)\}.
\]
反之，偶连续函数可写成
$f(x)=g(x^2)$，其中 $g(t)=f(\sqrt t)$ 连续于 $[0,1]$。由 Weierstrass
定理取多项式 $p_n\to g$，则偶多项式 $p_n(x^2)$ 一致趋于 $f$。故闭包
恰为 $E$。
\end{solution}

\begin{solution}{polynomials-with-zero-at-origin}
一致极限保持在零点取值为零，所以闭包包含于
$M=\{f\in C([0,1]):f(0)=0\}$。反之，对 $f\in M$，由 Weierstrass 取
$q_n\to f$，令 $p_n=q_n-q_n(0)$。则 $p_n(0)=0$，且
$\norm{p_n-f}\le\norm{q_n-f}+|q_n(0)-f(0)|\le2\norm{q_n-f}\to0$。
因此闭包就是 $M$。
\end{solution}

\begin{solution}{complex-conjugation-necessity}
在 $C(\mathbb T;\C)$ 上定义连续线性泛函
\[
 L(f)=\int_0^{2\pi} f(e^{it})e^{it}\,dt .
\]
由 $|L(f)|\le2\pi\norm f_\infty$ 可知 $L$ 对一致收敛连续。若
$p(z)=\sum_{k=0}^Na_kz^k$，则
\[
 L(p)=\sum_{k=0}^Na_k\int_0^{2\pi}e^{i(k+1)t}\,dt=0.
\]
因此 $L$ 在这类多项式的一致闭包上仍为零。然而
\[
 L(\bar z)=\int_0^{2\pi}e^{-it}e^{it}\,dt=2\pi.
\]
故 $\bar z$ 不属于该闭包。这一例说明：在复值版本中，含常数和分离点仍不
足以保证稠密，共轭封闭条件不能删去。
\end{solution}

\begin{solution}{bernstein-convexity-preservation}
当 $n=1$ 时，$B_1f$ 为仿射函数，故凸。对 $n\ge2$，二次求导整理为
\[
 (B_nf)''(x)=n(n-1)\sum_{k=0}^{n-2}
 \Delta^2f(k/n)\binom{n-2}kx^k(1-x)^{n-2-k},
\]
其中
$\Delta^2f(k/n)=f((k+2)/n)-2f((k+1)/n)+f(k/n)\ge0$ 对凸函数成立。
故 $B_nf$ 凸。对所有 $n\ge1$，有限凸组合的 Jensen 不等式给出
$f(x)=f(\sum_k(k/n)p_{n,k})\le\sum_kf(k/n)p_{n,k}=B_nf(x)$，所以
Bernstein 多项式从上方逼近凸函数。
\end{solution}
